\documentclass[11pt]{article}
\usepackage{fontspec}
\usepackage[french]{babel}
\usepackage{amsmath}
\usepackage{amssymb}
\usepackage{amsthm}
\usepackage{mathrsfs}
\usepackage{tikz-cd}
\usepackage{enumitem}

\newcommand{\printedpage}[1]{\marginpar{\footnotesize p.~#1}}
\newcommand{\msp}[1]{\marginpar{\footnotesize [#1]}}
\newcommand{\Dl}{\boldsymbol{\Delta}}
\newcommand{\tl}{\tilde{}}
\newcommand{\Ob}{\mathrm{Ob}}
\newcommand{\Fl}{\mathrm{Fl}}

\begin{document}

\printedpage{49}
% head: A. Grothendieck : Dérivateurs, Ch. XVII, Édition des Universités de Montpellier II et Paris VII
% foot: Version 9 février 2019

\msp{page 98}

\section*{6\quad Application à la construction des dérivateurs (associés aux triples de Quillen spéciaux)}

\textbf{Théorème 6.1.}\footnote[29]{C'est le résultat principal dans cette partie XVII des notes.} \textit{Soit $(W,C,F)$ un triple de Quillen clos spécial sur la catégorie $\mathcal{M}$} (cf. déf. 5.6, p. 90, et cor. 5.7). \textit{Pour tout $I$ dans} Cat\textit{, considérons la catégorie}
\[
\mathcal{M}(I) \;=\; \underline{\mathrm{Hom}}(I^{\circ},\mathcal{M}) \;, \tag*{(6.1.1)}
\]
\textit{et considérons}
\[
W(I),\ F(I) \;\subset\; \Fl(\mathcal{M}(I)) \;,
\]
\textit{définis par}
\[
\begin{aligned}
&\text{\textit{$u : X \longrightarrow Y$ dans $\mathcal{M}(I)$ est dans $W(I)$ (resp. $F(I)$) si et seulement}}\\
&\text{\textit{si $\forall\, i \in \Ob I$, $u(i) : X(i) \longrightarrow Y(i)$ est dans $W$ (resp. $F$).}}
\end{aligned}
\tag*{(6.1.2)}
\]

\begin{itemize}[leftmargin=2.5em]
\item[a)] \textit{Il existe une structure de Quillen close (nécessairement unique) dans $\mathcal{M}(I)$, ayant $W(I)$ et $F(I)$ comme ensembles des quasi-isomorphismes et de fibrations respectivement. On a donc aussi}
\[
\left\{
\begin{aligned}
TF(I) \;&=\; F(I)\cap W(I) \;=\; \{u\in \Fl(\mathcal{M}(I)) \mid u(i)\in TF \ \ \forall\, i\in \Ob I\},\\
C(I) \;&=\; TF(I)^{*}.
\end{aligned}
\right.
\tag*{(6.1.3)}
\]
\textit{Cette structure de Quillen est} spéciale.

\msp{page 99}
\item[b)] \textit{Soit $I_0$ la catégorie discrète ayant même ensemble d'objets que $I$, et}
\[
v_I \ \ \text{\textit{ou}}\ \ v \;:\; I_0 \longrightarrow I
\tag*{(6.1.4)}
\]
\textit{le foncteur canonique. Désignons par}
\[
C(I) \;=\; (TF(I))^{*}, \qquad TC(I) \;=\; C(I)\cap W(I) \;=\; F(I)^{*}
\tag*{(6.1.5)}
\]
\textit{les ensembles de cofibrations resp. de cofibrations triviales dans $\mathcal{M}(I)$, et utilisons les mêmes notations $C(I_0)$, $TC(I_0)$ pour $I_0$. On a alors}
\[
\begin{aligned}
&\text{\textit{$u : X_0 \longrightarrow Y_0$ dans $\mathcal{M}(I_0)$ est dans $C_0$ (resp. $TC_0$) si et seulement si}}\\
&\text{\textit{pour tout $i$ dans $\Ob I$, $u(i) : X_0(i) \longrightarrow Y_0(i)$ est dans $C$ (resp. $TC$).}}
\end{aligned}
\tag*{(6.1.6)}
\]
\textit{Enfin on a}
\[
\left\{
\begin{aligned}
C(I) \;&=\; (v_!(C(I_0)))\tl\\
TC(I) \;&=\; (v_!(TC(I_0)))\tl\,.
\end{aligned}
\right.
\tag*{(6.1.7)} % ?? uncertain: shape of the small superscript mark after the closing parenthesis -- tilde (\tilde{}) or breve (\breve{})
\]

\item[c)] \textit{Soit}
\[
u \;:\; I \longrightarrow I'
\]
\textit{une flèche dans} Cat\textit{, considérons le couple des foncteurs adjoints}
\[
\begin{tikzcd}
\mathcal{M}(I) \arrow[r, shift left, "u_!"] & \mathcal{M}(I') \arrow[l, shift left, "u^{*}"]
\end{tikzcd}
\;.
\tag*{(6.1.8)}
\]

\printedpage{50}
% head: A. Grothendieck : Dérivateurs, Ch. XVII, Édition des Universités de Montpellier II et Paris VII
% foot: Version 9 février 2019

\msp{page 100}
\textit{Ce couple satisfait aux conditions de Quillen pour les couples de foncteurs adjoints entre catégories de Quillen} (Hom. Alg. Chp. I, §4, th. 3, p. 4.5), \textit{de sorte qu'ils définissent par passage aux localisés un couple des foncteurs adjoints}
\[
\begin{tikzcd}
\mathrm{Ho}_W(I) \arrow[r, shift left, "u_!^{W}"] & \mathrm{Ho}_W(I') \arrow[l, shift left, "u^{*}_{W}"]
\end{tikzcd}
\;.
\tag*{(6.1.9)}
\]
\textit{On a, plus précisément, les relations}
\[
\left\{
\begin{aligned}
u_!(C(I)) \;&\subset\; C(I')\\
u_!(TC(I)) \;&\subset\; TC(I')\;,
\end{aligned}
\right.
\tag*{(6.1.10)}
\]
\textit{d'où}
\[
(u_!(\mathcal{M}(I))_{\mathrm{c}}) \;\subset\; \mathcal{M}(I')_{\mathrm{c}}\;,
\tag*{(6.1.11)}
\]
\textit{et de plus}
\[
\begin{aligned}
&\text{\textit{pour toute flèche $f$ dans $\mathcal{M}(I)_{\mathrm{c}}$, telle que $f \in W(I)$, on a}}\\
&\text{\textit{$u_!(f) \in W(I')$, i.e. $(u_!)_{\mathrm{c}} : \mathcal{M}(I)_{\mathrm{c}} \longrightarrow \mathcal{M}(I')_{\mathrm{c}}$ est compatible}}\\
&\text{\textit{avec les localiseurs.}}
\end{aligned}
\tag*{(6.1.12)}
\]
\textit{D'autre part on a}
\[
\left\{
\begin{aligned}
u^{*}(F(I')) \;&\subset\; F(I)\\
u^{*}(TF(I')) \;&\subset\; TF(I)\\
u^{*}(W(I')) \;&\subset\; W(I)\;,
\end{aligned}
\right.
\tag*{(6.1.13)}
\]
\textit{donc}
\[
\left\{
\begin{aligned}
&u^{*}(\mathcal{M}(I')_{\mathrm{f}}) \;\subset\; \mathcal{M}(I)_{\mathrm{f}}\\
&u^{*}_{\mathrm{f}} : \mathcal{M}(I')_{\mathrm{f}} \longrightarrow \mathcal{M}(I)_{\mathrm{f}} \text{ \textit{est compatible avec les localiseurs.}}
\end{aligned}
\right.
\tag*{(6.1.14)}
\]

\msp{page 101}
\item[d)] \textit{Supposons que $I$ soit une catégorie ordonnée, et considérons l'ensemble $C_I \subset \Fl(\mathcal{M}(I))$ défini dans §2}\footnote[30]{Au paragraphe 2, c'était noté $C_{I^{\circ}}$ au lieu de $C_I$.} (cf. page 25, cf. aussi th. 2.1, p. 26, et ses corollaires). \textit{On a alors}
\[
\left\{
\begin{aligned}
C_I \;&\supset\; C(I)\\
TC_I \;&\supset\; TC(I)\;,
\end{aligned}
\right.
\tag*{(6.1.15)}
\]
\textit{et on a égalité quand $I$ est un ensemble ordonné} artinien.
\end{itemize}

\textsc{Démonstration}. L'assertion a) est un cas particulier de 5.8, appliqué au couple de foncteurs adjoints
\[
\begin{tikzcd}
\mathcal{M}(I_0) \arrow[r, shift left, "v_!"] & \mathcal{M}(I) \arrow[l, shift left, "v^{*}"]
\end{tikzcd}
\;.
\]
Les conditions (5.8.2) sont satisfaites, cela a été vu dans (4.8). La proposition 5.8 nous dit alors l'existence d'une structure de triple de Quillen close sur $\mathcal{M}(I)$, dont l'ensemble des quasi-isomorphismes $W(I)$ et l'ensemble des fibrations $F(I)$, sont donnés par la première formule (5.8.3) et par la première formule (5.8.4), ce qui équivaut aux définitions (6.1.2). (\textbf{N.B.} Le fait que $\mathcal{M}(I_0)$, avec la structure $W(I_0)$, $C(I_0)$, $F(I_0)$ évidente (produit), est
\printedpage{51}
% head: A. Grothendieck : Dérivateurs, Ch. XVII, Édition des Universités de Montpellier II et Paris VII
% foot: Version 9 février 2019
\msp{page 102}
bien une catégorie de Quillen spéciale close, est $\pm$ tautologique.)

\textit{Prouvons} b). La relation (6.1.6) est tautologique. Les relations (6.1.7) sont contenues dans 5.8 (cf. deuxième relation dans (5.8.3) pour $C' = v_!(C)\tl$, i.e. $C(I) = (v_!(C(I_0)))\tl$, et (5.8.5) pour $TC(I) = (v_!(TC(I_0)))\tl$).

\textit{Prouvons} c). La première assertion est, par \textsc{Quillen}, un cas particulier des relations (6.1.10) à (6.1.14), c'est elles qu'il s'agit d'établir. Les relations (6.1.13) et par suite (6.1.14) sont tautologiques en vertu des définitions, donc il s'agit d'établir \textit{(6.1.10)} et (6.1.12).

Pour les formules (6.1.10), j'utilise le diagramme commutatif
\[
\begin{tikzcd}[row sep=large, column sep=large]
I \arrow[d, "u"'] & I_0 \arrow[l, "v"'] \arrow[d, "u_0"] \\
I' & I'_0\;, \arrow[l, "v'"']
\end{tikzcd}
\]
et je note que les foncteurs $u_!^{\mathcal{M}}$, $u_{0\,!}^{\mathcal{M}}$, commutant aux petites limites inductives, se
\msp{page 103}
comportent bien pour l'opération
\[
\Phi \longmapsto \Phi\tl\,,
\]
quand $\Phi$ est une partie de $\Fl(\mathcal{M}(I))$ resp. de $\Fl(\mathcal{M}(I_0))$ qui est Q-accessible, i.e. telle qu'il existe une \textit{petite} partie $\Phi_0 \subset \Phi$ avec $\Phi \subset \tilde{\Phi}_0$, plus précisément, on a
\[
u_!(\Phi\tl) \;\subset\; (u_!(\Phi))\tl\,.
\]
Ceci noté, on a par b)
\[
C(I) \;=\; (v_!(C(I_0)))\tl\,,
\]
d'où
\[
u_!(C(I)) \;\subset\; (\underbrace{u_!(v_!(C(I_0)))}_{(uv)_!(C(I_0))})\tl\,,
\]
et le deuxième membre s'écrit aussi
\[
(v'_!u_{0\,!}(C(I_0)))\tl \;\overset{?}{\subset}\; (v'_!(I'_0))\tl \;=\; C(I')\;. % [sic] (v'_!(I'_0)) -> (v'_!(C(I'_0))), the C( ) around I'_0 apparently omitted
\]
Donc tout revient à montrer que
\[
u_{0\,!}(C(I_0)) \;\subset\; C(I'_0)\;,
\tag*{(6.1.16)}
\]
ce qui résulte aussitôt du calcul de $u_{0\,!}$ à coups de sommes directes, et du fait que $C$ est stable par sommes directes quelconques. Cela prouve donc $u_!(C(I)) \subset C(I')$, et on prouve de même $u_!(TC(I)) \subset TC(I')$, qui se ramène de même au cas particulier
\[
u_{0\,!}(TC(I_0)) \;\subset\; TC(I'_0)\;.
\tag*{(6.1.17)}
\]

\printedpage{52}
% head: A. Grothendieck : Dérivateurs, Ch. XVII, Édition des Universités de Montpellier II et Paris VII
% foot: Version 9 février 2019

\msp{page 104}

Il reste à prouver (6.1.12), et on procède comme pour 3.3 (p. 54), en utilisant le fait qu'une flèche $f$ dans $\mathcal{M}(I)_{\mathrm{c}}$ qui est dans $W(I)$, s'écrit dans $\mathcal{M}(I)_{\mathrm{c}}$ sous la forme $f = pi$, où $i \in TC(I)$, et où $p$ est inverse à gauche (rétraction) d'une flèche $j \in TC(I)$. On aura donc $u_!(f) = u_!(p)u_!(i)$, et $u_!(i) \in TC(I')$ par (6.1.10), et $u_!(p)$ est inverse à gauche de $u_!(j) \in TC(I')$. On a alors $u_!(i), u_!(j), u_!(p) \in W(I')$, d'où $u_!(f) \in W(I')$, q.e.d.

Prouvons enfin d). Compte tenu de la définition de $C_I$, pour établir $C(I) \subset C_I$, il suffit de voir ceci :

\textbf{Corollaire 6.2.}\footnote[31]{\textbf{N.B.} Ce résultat se généralise, sous les conditions générales dans 6.2.5.} \textit{Soit $I$ quelconque dans} Cat\textit{, soient $U \subset V$ deux fermés dans $I$ (i.e. des ouverts dans $I^{\circ}$), et soit $f : X_{\bullet} \longrightarrow Y_{\bullet}$ une flèche dans $C(I)$ (resp. dans $TC(I)$). Posons encore}
\msp{page 105}
\[
X_U \;=\; \varinjlim_{i\in U} X(i), \qquad X_V \;=\; \varinjlim_{i\in U} X(i)\;, % [sic] second index i\in U -> i\in V, X_V should be the limit over V
\]
\textit{et de même pour $Y$, et considérons le diagramme}
\[
\begin{tikzcd}[row sep=large, column sep=large]
X_U \arrow[r, "\text{trans.}"] \arrow[d, "u_U"'] \arrow[dr, phantom, "\text{cocart.}"] & X_V \arrow[d] \arrow[dd, dotted, bend left=60, "u_V"] \\
Y_U \arrow[r] \arrow[dr, dotted, "\text{trans.}"'] & Y_{U,V} \arrow[d, "u_{U,V}"'] \\
 & Y_V\;.
\end{tikzcd}
\tag*{(6.2.1)}
\]
\textit{Alors la flèche $u_{U,V}$ est dans $C$ (resp. $TC$).}

\textsc{Démonstration} de 6.2. Admettons que la restriction d'une flèche de $C(I)$ resp. $TC(I)$ à un fermé $V$ de $I$ est dans $C(V)$ resp. $TC(V)$ (cf. \dots{} ci-dessous). Cela nous permet de supposer $V = I$. La donnée de $U$, qui est un ouvert de $I^{\circ}$, équivaut à la donnée d'une flèche $v : I^{\circ} \longrightarrow \Dl^1$ (en prenant la fibre en 0), ou encore à la donnée d'une flèche $u : I \longrightarrow (\Dl^1)^{\circ} \simeq \Dl^1$ (en prenant le fibre en 1). Or % [sic] le fibre -> la fibre, gender (cf. "la fibre en 0" just before)
\[
u \;:\; I \longrightarrow \Dl^1\;, \qquad U \;=\; I_1 \overset{\text{déf}}{=} u^{-1}(\{1\})\;.
\tag*{(6.2.2)}
\]

\msp{page 106}
On a alors
\[
X_U \;=\; u_!(X)(1), \quad X_I \;=\; u_!(X)(0)\;,
\]
et le morphisme de transition $X_U \longrightarrow X_I$ est déduit de la flèche $0 \xrightarrow{\ \partial\ } 1$ dans $\Dl^1$, et de même pour $Y$. Soit
\[
\xi \;=\; u_!(X), \qquad \eta \;=\; u_!(Y)\;,
\]
\printedpage{53}
% head: A. Grothendieck : Dérivateurs, Ch. XVII, Édition des Universités de Montpellier II et Paris VII
% foot: Version 9 février 2019
donc le diagramme (6.2.1) s'identifie au diagramme déduit de
\[
\varphi \;:\; \xi \longrightarrow \eta\;, \qquad \varphi \;=\; u_!(f)\;,
\]
comme
\[
\begin{tikzcd}[row sep=large, column sep=large]
\xi(1) \arrow[r, "\xi(\partial)"] \arrow[d, "\varphi(1)"'] \arrow[dr, phantom, "\text{cocart.}"] & \xi(0) \arrow[d] \arrow[dd, dotted, bend left=60, "\varphi(0)"] \\
\eta(1) \arrow[r] \arrow[dr, dotted, "\eta(\partial)"'] & \tilde{\eta} \arrow[d, "\tilde{\varphi}"'] \\
 & \eta(0) \qquad .
\end{tikzcd}
\tag*{(6.2.3)}
\]
Or par 6.1.10 appliqué à $u : I \longrightarrow \Dl^1$, la relation $f \in C(I)$ resp. $f \in TC(I)$ implique $\varphi \in C(\Dl^1)$ resp. $\varphi \in TC(\Dl^1)$. Donc pour établir 6.1, on est ramené au cas particulier où $I = \Dl^1$, $U = \{1\}$, $V = I$, donc à prouver le

\textbf{Lemme 6.2.4.} \textit{Soit $\varphi : \xi \longrightarrow \eta$ une flèche dans $\mathcal{M}(\Dl^1)$. Si $\varphi \in C(\Dl^1)$ (resp.}
\msp{page 107}
\textit{$\varphi \in TC(\Dl^1)$), alors la flèche $\tilde{\varphi}$ dans le diagramme} (6.2.3) \textit{est dans $C$ (resp. $TC$).}

Plus généralement, et plus précisément :

\textbf{Lemme 6.2.5.} \textit{Soit $\mathcal{M}$ une catégorie stable par sommes amalgamées, et $\Phi$ une partie Q-close à gauche} (on prendra $\Phi = C$, et $\Phi = TC$, dans le cas du lemme 6.2.4). \textit{Soit $\Psi = \Phi_{*}$, et considérons, pour toute $I$ dans} Cat\textit{, $\Phi(I)$, $\Psi\langle I\rangle$ dans $\mathcal{M}(I)$, définis par :}
\[
\begin{aligned}
&\text{\textit{$u : X \longrightarrow Y$ dans $\mathcal{M}(I)$ est dans $\Psi\langle I\rangle$ si et seulement si $\forall\, i \in \Ob I$,}}\\
&\text{\textit{$u(i) : X(i) \longrightarrow Y(i)$ est dans $\Psi$,}}
\end{aligned}
\tag*{(6.2.5.1)}
\]
\[
\Phi(I) \;=\; \Psi\langle I\rangle^{*}.
\tag*{(6.2.5.2)}
\]

Prenons le cas $I = \Dl^1 = \{0 \xrightarrow{\ \partial\ } 1\}$, et soit $u : X \longrightarrow Y$ une flèche dans $\mathcal{M}(\Dl^1)$, correspondant à un carré commutatif
\[
\begin{tikzcd}[row sep=large, column sep=large]
X(1) \arrow[r, "X(\partial)"] \arrow[d, "u(1)"'] & X(0) \arrow[d, "u(0)"] \\
Y(1) \arrow[r, "Y(\partial)"] & Y(0)\;.
\end{tikzcd}
\tag*{(6.2.5.3)}
\]

Considérons le diagramme correspondant

\end{document}
